\documentclass[11pt]{article}
\usepackage[margin=1in]{geometry}
\usepackage[T1]{fontenc}
\usepackage[utf8]{inputenc}
\usepackage{lmodern}
\usepackage{setspace}
\usepackage{amsmath, amssymb}
\usepackage{booktabs}
\usepackage{enumitem}
\usepackage{xcolor}
\usepackage{hyperref}
\usepackage{titlesec}

\definecolor{titleblue}{RGB}{23,54,93}
\hypersetup{
    colorlinks=true,
    linkcolor=titleblue,
    urlcolor=titleblue,
    citecolor=titleblue
}

\titleformat{\section}{\large\bfseries\color{titleblue}}{\thesection.}{0.5em}{}
\titleformat{\subsection}{\normalsize\bfseries\color{titleblue}}{\thesubsection}{0.5em}{}
\setlist[itemize]{leftmargin=1.5em}
\setstretch{1.12}

\title{\textbf{Learning-Augmented Robust Traffic Engineering for Resilient Software-Defined Networks}\\[0.3em]
\large Final Project Proposal}
\author{ECE GY 7363: Communications Networks II -- Design and Algorithms}
\date{April 2026}

\begin{document}
\maketitle

\begin{abstract}
This project proposes a learning-augmented framework for traffic engineering in software-defined networks under time-varying demand and link-failure uncertainty. The main idea is to combine short-horizon traffic prediction with optimization-based multi-commodity flow routing so that the controller can proactively reduce congestion while preserving performance under failures. The machine learning component predicts the next traffic matrix from historical observations, while the optimization component computes routing decisions that minimize maximum link utilization subject to flow conservation and capacity constraints. The resulting framework will be evaluated against shortest-path routing, optimization without prediction, and conservative robust baselines on standard network topologies and synthetic dynamic traffic traces. The proposal directly addresses course themes in network design modeling, optimization, fairness, resilient design, robust design, and machine-learning-based network management. The expected outcome is a systematic study of when prediction-enhanced traffic engineering improves efficiency and resilience, and when robust non-predictive methods remain preferable.
\end{abstract}

\section{Introduction and Motivation}
Communication networks increasingly operate in environments where traffic demand changes rapidly over time and failures can occur with little warning. Traditional shortest-path routing is simple and computationally cheap, but it often reacts poorly to congestion because it does not explicitly reason about capacity sharing across commodities. Optimization-based traffic engineering improves this situation by accounting for the full demand matrix and link capacities, yet a purely reactive optimization approach may still perform suboptimally when demand shifts between control intervals. At the same time, machine learning has become a practical tool for traffic forecasting, anomaly detection, and resource management, but ML alone does not guarantee feasible routing decisions or provable capacity compliance.

This project is motivated by the observation that prediction and optimization are complementary. If an accurate predictor can estimate the near-future traffic matrix, then a routing optimizer may use that information to allocate capacity more effectively before congestion materializes. However, prediction errors and network failures introduce uncertainty, so the design must remain robust even when forecasts are imperfect. This leads to the central idea of the proposal: a learning-augmented traffic engineering framework that combines demand prediction, multi-commodity flow optimization, and resilience analysis in a single end-to-end network design study.

\section{Problem Statement}
The project considers a communication network represented by a directed graph
\[
G = (V, E),
\]
where $V$ is the set of nodes and $E$ is the set of directed links. Each link $(i,j) \in E$ has capacity $c_{ij} > 0$. At each decision epoch $t$, the network is presented with a traffic matrix containing demands between source-destination pairs. Each source-destination pair is treated as a commodity $k \in \mathcal{K}$ with demand $d_k(t)$.

The design problem is to compute routing decisions that satisfy as much demand as possible while controlling congestion and maintaining good performance under failures. Specifically, the project asks the following research question:

\begin{quote}
Can short-term traffic prediction improve routing quality and resilience relative to shortest-path and non-predictive optimization baselines in dynamic software-defined networks?
\end{quote}

The project will study this question under traffic variability, traffic spikes, and link-failure scenarios.

\section{Project Objectives}
The proposal has five concrete objectives:
\begin{itemize}
    \item Formulate a dynamic traffic engineering model for a capacitated communication network.
    \item Train a machine learning model to predict near-future traffic matrices from historical data.
    \item Integrate predicted demand into a multi-commodity flow routing optimizer.
    \item Evaluate resilience under random and critical link-failure scenarios.
    \item Compare the proposed method against multiple classical and optimization-based baselines.
\end{itemize}

\section{Relevance to Course Themes}
The project is strongly aligned with the course description and final project expectations. It directly addresses:
\begin{itemize}
    \item \textbf{Network Design Problem Modeling:} formalizing dynamic traffic engineering as a capacitated graph optimization problem.
    \item \textbf{Optimization Methods:} solving linear or convex relaxations of multi-commodity flow routing.
    \item \textbf{Multi-Commodity Flow Routing:} explicitly routing concurrent source-destination demands over shared infrastructure.
    \item \textbf{Fair Network:} evaluating or extending the formulation with fairness metrics such as Jain's index.
    \item \textbf{Resilient and Robust Network Design:} testing routing performance under failures and uncertain prediction quality.
    \item \textbf{Machine Learning Based Network Management:} using sequence models to guide routing decisions.
\end{itemize}

\section{Technical Approach}
\subsection{System Overview}
At each time step, the controller will observe recent traffic matrices, use a predictive model to estimate the next traffic matrix, solve a routing optimization problem using the predicted demand, and then evaluate the selected routing under the realized traffic and failure scenarios. The resulting pipeline is:
\[
\text{historical traffic} \rightarrow \text{prediction model} \rightarrow \text{routing optimizer} \rightarrow \text{performance evaluation}.
\]

\subsection{Machine Learning Component}
The ML module will take a sliding window of recent traffic matrices as input and output a one-step-ahead traffic matrix prediction. The study will include:
\begin{itemize}
    \item a simple baseline predictor such as moving average or linear regression,
    \item a feed-forward multilayer perceptron baseline if useful,
    \item a recurrent model such as LSTM or GRU as the main predictive approach.
\end{itemize}

The primary model will likely be an LSTM because dynamic traffic has temporal correlation and burst structure that sequence models can capture effectively. Prediction performance will be measured using mean absolute error, root mean squared error, and relative error on traffic entries and aggregate demand.

\subsection{Optimization Component}
The routing problem will be modeled as a multi-commodity flow program. Let $f_{ij}^{k}(t)$ denote the amount of commodity $k$ routed on link $(i,j)$ at decision epoch $t$. A standard formulation uses the following constraints:

\paragraph{Flow conservation.}
For each commodity $k$ and node $v \in V$,
\[
\sum_{(u,v)\in E} f_{uv}^{k}(t) - \sum_{(v,w)\in E} f_{vw}^{k}(t) =
\begin{cases}
-d_k(t), & v = s_k, \\
d_k(t), & v = t_k, \\
0, & \text{otherwise,}
\end{cases}
\]
where $s_k$ and $t_k$ are the source and destination of commodity $k$.

\paragraph{Capacity constraints.}
For every link $(i,j)\in E$,
\[
\sum_{k\in\mathcal{K}} f_{ij}^{k}(t) \le c_{ij}.
\]

\paragraph{Nonnegativity.}
\[
f_{ij}^{k}(t) \ge 0 \quad \forall (i,j)\in E, \; k\in\mathcal{K}.
\]

To minimize congestion, the project will introduce an auxiliary utilization variable $U(t)$ and solve:
\[
\min \; U(t)
\]
subject to
\[
\sum_{k\in\mathcal{K}} f_{ij}^{k}(t) \le U(t) \, c_{ij}, \quad \forall (i,j)\in E.
\]
This yields a minimum-max-utilization routing policy, a standard and interpretable objective in traffic engineering.

\subsection{Robustness and Resilience}
The project will study resilience in two layers:
\begin{itemize}
    \item \textbf{Scenario-based resilience evaluation:} routing is optimized under nominal conditions and then tested under link-failure scenarios.
    \item \textbf{Optional robust extension:} if time allows, a simplified robust optimization model will be developed to hedge against a prespecified set of failure scenarios or demand perturbations.
\end{itemize}

The first layer ensures that the project always contains a clear resilient-network component even if a full robust formulation is computationally expensive.

\section{Algorithms and Baselines}
The proposed method will be compared against the following baselines:
\begin{enumerate}
    \item \textbf{Shortest-path routing:} each demand is routed using a fixed path metric without optimization or prediction.
    \item \textbf{Optimization with current demand only:} multi-commodity flow routing is computed using the currently observed traffic matrix.
    \item \textbf{Conservative non-predictive routing:} a robust or safety-margin-based formulation uses inflated demand estimates without ML.
    \item \textbf{Learning-augmented routing (proposed):} the predictor estimates the next traffic matrix, and the optimization uses the predicted demand to compute routes.
\end{enumerate}

This comparison will isolate the value of ML, the value of optimization, and the value of combining both.

\section{Data, Topologies, and Experimental Setup}
\subsection{Network Topologies}
The experiments will use one or more standard topologies such as NSFNET, Abilene, or GEANT. If necessary, additional synthetic topologies will be generated to examine stress cases or larger-scale settings.

\subsection{Traffic Generation}
The study will use time-varying traffic matrices created from a combination of:
\begin{itemize}
    \item periodic demand patterns,
    \item random fluctuations,
    \item bursty spikes,
    \item hotspot source-destination pairs.
\end{itemize}
This provides realistic temporal structure while keeping the dataset generation reproducible and fully documented.

\subsection{Failure Scenarios}
The resilience study will include:
\begin{itemize}
    \item no-failure baseline cases,
    \item random single-link failures,
    \item critical-link failures selected by betweenness or capacity centrality,
    \item traffic-spike cases combined with failures.
\end{itemize}

\section{Evaluation Metrics}
The following metrics will be reported:
\begin{itemize}
    \item maximum link utilization,
    \item average link utilization,
    \item fraction of demand satisfied,
    \item total routing cost,
    \item average path length,
    \item number of congested links,
    \item Jain's fairness index,
    \item degradation under failures,
    \item prediction error of the traffic forecasting model.
\end{itemize}

The principal metric will be maximum link utilization because it directly captures congestion severity and is central to traffic engineering performance.

\section{Expected Contributions}
The anticipated contributions of the project are:
\begin{itemize}
    \item a complete learning-plus-optimization traffic engineering framework,
    \item a rigorous comparison of predictive and non-predictive routing methods,
    \item a resilience analysis under traffic uncertainty and failures,
    \item practical insight into the operating regimes where ML improves network design decisions.
\end{itemize}

Even if the ML model does not dominate in every setting, identifying when prediction helps and when conservative optimization is preferable is itself a meaningful result.

\section{Implementation Plan}
The project will be implemented in Python. The expected software stack is summarized in Table~\ref{tab:stack}.

\begin{table}[h]
\centering
\begin{tabular}{ll}
\toprule
\textbf{Component} & \textbf{Planned Tooling} \\
\midrule
Graph and topology processing & \texttt{networkx} \\
Numerical processing & \texttt{numpy}, \texttt{pandas} \\
Optimization & \texttt{cvxpy} or \texttt{pulp}; optional Gurobi/CPLEX \\
Machine learning & \texttt{PyTorch} \\
Visualization & \texttt{matplotlib}, \texttt{seaborn} \\
\bottomrule
\end{tabular}
\caption{Planned implementation stack.}
\label{tab:stack}
\end{table}

The implementation will be organized into modules for topology loading, traffic generation, forecasting, optimization, failure simulation, and result visualization. This modular structure will make experiments easier to reproduce and present.

\section{Work Plan and Timeline}
The project will proceed in the following sequence:
\begin{enumerate}
    \item Finalize topology selection, traffic model, and evaluation metrics.
    \item Implement shortest-path routing and basic evaluation code.
    \item Formulate and validate the multi-commodity flow optimization model.
    \item Train baseline and LSTM traffic prediction models.
    \item Integrate the predicted traffic matrix with the optimizer.
    \item Run experiments under dynamic traffic and link failures.
    \item Analyze results and prepare the final report, code package, and presentation.
\end{enumerate}

A tentative week-by-week schedule is shown below:
\begin{itemize}
    \item \textbf{Week 1:} problem setup and codebase organization.
    \item \textbf{Week 2:} synthetic traffic generation and shortest-path baseline.
    \item \textbf{Week 3:} optimization model implementation.
    \item \textbf{Week 4:} ML model training and validation.
    \item \textbf{Week 5:} integration of prediction and routing.
    \item \textbf{Week 6:} resilience experiments and sensitivity analysis.
    \item \textbf{Week 7:} documentation, slides, and final packaging.
\end{itemize}

\section{Risks and Mitigation}
There are three main risks. First, the optimization model may become slow on larger topologies. This will be managed by beginning with medium-size networks and keeping a simpler baseline available. Second, traffic prediction may not produce meaningful improvements if the generated traffic is too noisy. This will be mitigated by constructing traffic traces with controlled temporal structure and by evaluating prediction accuracy independently from routing performance. Third, a full robust optimization model may be too costly to finish within the available time. The project therefore prioritizes scenario-based resilience evaluation first and treats the robust formulation as an extension.

\section{Why the Project is Well Positioned for a High Score}
The proposed study is well aligned with the final project rubric. It is relevant to the course material, technically deep, grounded in mathematical optimization, and strengthened by a modern ML component. The results can be presented clearly through traffic matrices, network diagrams, utilization plots, failure-case comparisons, and fairness metrics. Most importantly, the project is ambitious without being unrealistic: even a reduced version still contains a coherent model, a meaningful algorithmic contribution, and a rigorous experimental design.

\section{Conclusion}
This proposal presents an end-to-end study of learning-augmented traffic engineering for resilient software-defined networks. By combining short-term traffic prediction with optimization-based multi-commodity flow routing, the project will investigate whether forecast-informed decisions can reduce congestion and improve robustness under failures. The proposed framework is mathematically grounded, experimentally testable, and directly aligned with the goals of the course. As a result, it offers a strong balance of novelty, depth, correctness, and presentation quality for the final design project.

\end{document}
